Small code generation models exhibit high syntax error rates that render outputs unusable before semantic evaluation---CodeLlama-7B fails to generate parseable Python for 71\% of HumanEval problems. We introduce syntax-aware beam search, which treats syntax validity as an explicit scoring dimension ($\beta=0.3$ weight) combined with log-likelihood ($\alpha=0.7$) to guide generation toward syntactically correct outputs. Unlike grammar-based constrained decoding that enforces hard constraints at high computational cost, type-constrained decoding that targets minority failure modes with weak penalties, or pure beam search that ignores syntax entirely, our approach occupies a middle ground: beam search ($k=5$) explores multiple candidate paths while lightweight AST validation ($<0.05$ms per check) provides sufficient signal to systematically prune invalid alternatives. We validate the complete pipeline through five mechanistic hypotheses testing infrastructure feasibility, beam diversity, scoring effectiveness, pruning dynamics, and final output quality, demonstrating that each component functions as designed with large safety margins (12-40 percentage points above targets). Results on HumanEval with CodeLlama-7B show 76.22\% final validity (23.78\% error rate), representing 66.4\% relative error reduction from the 70.73\% baseline, at 4.5$\times$ computational cost versus greedy sampling (15 minutes for 164 problems). This 3.2$\times$ improvement in usable outputs makes small models viable for resource-constrained applications---IDE autocomplete, educational tools, internal automation---where deployment costs or latency requirements prevent using expensive large models.
